\clearpage
\section*{Appendix R\quad Computational Receipts}\label{app:receipts}
\addcontentsline{toc}{section}{Appendix R: Computational Receipts}
\small\noindent Each computation marked \rcptmarker\ in the text is verified by a runnable receipt in the \texttt{receipts/} directory that ships with this work; status \textsf{OK} = verified (traced, run, output matches the claim). This appendix is generated from the receipt index, not maintained by hand.\par\medskip
\begingroup\renewcommand{\arraystretch}{1.25}
\begin{description}
\item[\label{rcpt:P02_cycloid_and_critical_points}\texttt{P02\_cycloid\_and\_critical\_points}]\hfill\textsf{[OK]}\\
\textit{prop:cycloid/prop:critical} \ (cycloid r=M(1+cos \ensuremath{\eta}), \ensuremath{\tau}=M(\ensuremath{\eta}+sin \ensuremath{\eta}) solves the interior radial eq; 2 critical points = r-poles of circle (r\ensuremath{-}M)\textsuperscript{2}+s\textsuperscript{2}=M\textsuperscript{2}). 
\emph{Computes:} substitutes cycloid into d\ensuremath{\tau}\textsuperscript{2}=r/(2M\ensuremath{-}r)dr\textsuperscript{2} + differentiates for the non-degenerate critical points; control: wrong \ensuremath{\tau} fails
 \emph{Bound:} parametrisation+critical points only (continuations/Kretschmann/ring separate)
\ \emph{Run:} \texttt{python3 receipts/P02\_janzen\_circle/P02\_cycloid\_and\_critical\_points.py}
\item[\label{rcpt:P02_bead_K_mass_free}\texttt{P02\_bead\_K\_mass\_free}]\hfill\textsf{[OK]}\\
\textit{sec:kretschmann (bead reading)} \ (K is MASS-FREE along the cosmological branch (64/27 s\textasciicircum{}4) and goes as M\textasciicircum{}-4 along the interior cycloid: the mass-dependence tracks the parametrisation being read, not the geometry read through it). 
\emph{Computes:} symbolic K on both readings; amplitude r \textasciitilde{} M\textasciicircum{}(1/3) gives r\textasciicircum{}6 \textasciitilde{} M\textasciicircum{}2 so M\textasciicircum{}2/r\textasciicircum{}6 \textasciitilde{} M\textasciicircum{}0 (6 x 1/3 = 2 exactly); cycloid amplitude r \textasciitilde{} M gives M\textasciicircum{}-4
 \emph{Bound:} the ONTOLOGICAL reading of the residual divergence is left open by the paper; this is the parametrisation fact
\ \emph{Run:} \texttt{python3 receipts/P02\_janzen\_circle/P02\_bead\_K\_mass\_free.py}
\item[\label{rcpt:P02_kretschmann_chain_rule}\texttt{P02\_kretschmann\_chain\_rule}]\hfill\textsf{[OK]}\\
\textit{sec:kretschmann} \ (K-divergence at r=0 is a chart-labelling/chain-rule artefact (12th-order pole; label swap moves it to the horizon)). 
\emph{Computes:} K(z)=48/(M\textsuperscript{4}(1+cos z)\textsuperscript{6}), Laurent pole order exactly 12 = 2\ensuremath{\times}6, relabel r=M(1\ensuremath{-}cos z) moves pole to z=0
 \emph{Bound:} label-tracking fact, not the ontological reading (left open by paper)
\ \emph{Run:} \texttt{python3 receipts/P02\_janzen\_circle/P02\_kretschmann\_chain\_rule.py}
\item[\label{rcpt:P02_analytic_continuations}\texttt{P02\_analytic\_continuations}]\hfill\textsf{[OK]}\\
\textit{sec:continuation} \ (z=i\ensuremath{\rho} recovers Schwarzschild exterior (Region I); z=\ensuremath{\pi}+i\ensuremath{\rho}' continues onto r<0 (back-seam)). 
\emph{Computes:} trig\ensuremath{\to}hyperbolic subs; 1\ensuremath{-}2M/r=tanh\textsuperscript{2}(\ensuremath{\rho}/2), radial+angular coeffs match; r=M(1\ensuremath{-}cosh \ensuremath{\rho}')<0 back-seam
 \emph{Bound:} metric identities of the two continuations; global Kruskal structure separate
\ \emph{Run:} \texttt{python3 receipts/P02\_janzen\_circle/P02\_analytic\_continuations.py}
\item[\label{rcpt:P02_ring_lambda_limit}\texttt{P02\_ring\_lambda\_limit}]\hfill\textsf{[OK]}\\
\textit{sec:ring} \ (single horizon = \ensuremath{\Lambda}\ensuremath{\to}0 limit of SdS root triple (2 roots \ensuremath{\to}\ensuremath{\infty}, 1 \ensuremath{\to}2M; triple sums to 0)). 
\emph{Computes:} cubic r\textsuperscript{3}\ensuremath{-}\ensuremath{\alpha}\textsuperscript{2}r+2M\ensuremath{\alpha}\textsuperscript{2}=0, Vieta sum=0, \ensuremath{\alpha}-scan finite root\ensuremath{\to}2M others\ensuremath{\to}\ensuremath{\pm}\ensuremath{\infty}; control: finite \ensuremath{\alpha}\ensuremath{\to}3 horizons
 \emph{Bound:} horizon-count only; A\ensuremath{_2}/generation readings cited not established
\ \emph{Run:} \texttt{python3 receipts/P02\_janzen\_circle/P02\_ring\_lambda\_limit.py}
\item[\label{rcpt:P02_interior_metric}\texttt{P02\_interior\_metric}]\hfill\textsf{[OK]}\\
\textit{lem:interior\_metric} \ (DERIVES the Schwarzschild interior metric in cycloid coords (eta,t,theta,phi): g\_tt=tan\textasciicircum{}2(eta/2) with r=M(1+cos eta), eta timelike / t spacelike, cycloid radial eqn g\_rr dr\textasciicircum{}2=-dtau\textasciicircum{}2). 
\emph{Computes:} derives g\_tt,g\_rr,angular from Schwarzschild components; signature check
\ \emph{Run:} \texttt{python3 receipts/P02\_janzen\_circle/P02\_interior\_metric.py}
\item[\label{rcpt:P02_the_third_axis_is_two_poles}\texttt{P02\_the\_third\_axis\_is\_two\_poles}]\hfill\textsf{[OK]}\\
\textit{sec:421 (the third complex axis)} \ (**THE COMPLEX-z CONTINUATIONS CLASSIFIED EXACTLY, AND THE COUNT OF THEM READ OFF THE CIRCLE.** ** THE REAL-LANDING SET IS EXACTLY THREE AXES AND THERE IS NO FOURTH: ** with z = a + ib the cycloid r = M(1+cos z) has **Im r = -M sin(a) sinh(b) -- a PRODUCT** -- so Im r = 0 iff b = 0 or a in pi*Z, and 2pi-periodicity collapses the second to a = 0 and a = pi. ** AND THE THREE EXHAUST THE GEOMETRY rather than merely the reality condition **: their ranges 0 <= r <= 2M (cycloid arc), r >= 2M (front seam) and r <= 0 (back seam) PARTITION the real radius line. On them tau is real, purely imaginary, and purely imaginary about tau = M pi -- **Lorentzian on the arc and Riemannian on the two seams**, which is what makes the imaginary axes physical rather than formal. ** AND THE COUNT IS THE CIRCLE'S OWN: ** the parametrisation carries z -> -z, under which **r is EVEN and tau is ODD**; its fixed points on R/2piZ are z = 0 and z = pi, where **r = 2M and r = 0 -- the event horizon and the curvature singularity, the paper's two poles**. So the number of continuations beyond the real one is **\textbar{}Fix(z -> -z)\textbar{} = 2**, not a fact about the cosine.). 
\emph{Computes:} Im r asserted equal to -M sin a sinh b; the r-ranges on both imaginary axes asserted; Re(tau) asserted zero on one and M*pi on the other; r asserted even and tau odd under z -> -z; r(0) = 2M and r(pi) = 0 asserted
 \emph{Bound:} ** THIS PROVES WHAT THE PAPER ASSERTED AND ADDS THE COUNT; IT OVERTURNS NOTHING. ** P2 already states the criterion, the list and the verdict. Not settled here and not needed for the count: whether the reflection z -> -z (r-even, tau-odd, i.e. cycloid-time reversal) is the corpus's R, its K, or a third thing
\ \emph{Run:} \texttt{python3 receipts/P02\_janzen\_circle/P02\_the\_third\_axis\_is\_two\_poles.py}
\item[\label{rcpt:P02_the_approach_is_mass_free}\texttt{P02\_the\_approach\_is\_mass\_free}]\hfill\textsf{[OK]}\\
\textit{sec:singularities (the two species)} \ (**THE SECOND AXIS THE METRIC-SINGULARITY CLASSIFICATION DOES NOT CARRY, and the negative answer it gives to the WORLDLINE half of the branch-point question.** ** Read at fixed r the divergence carries the mass, K = 48 M\textasciicircum{}2/r\textasciicircum{}6 + 24/alpha\textasciicircum{}4. Read along the marginally-bound infalling geodesic as a function of the FALLER'S OWN PROPER TIME, it does not: ** on r(tau) = (2M alpha\textasciicircum{}2)\textasciicircum{}(1/3) sinh\textasciicircum{}(2/3)(3 tau/2 alpha) the cube of the amplitude cancels the M\textasciicircum{}2 exactly, leaving **K = 12[2 + sinh\textasciicircum{}-4(3 tau/2 alpha)]/alpha\textasciicircum{}4 with no M anywhere**, and K -> (64/27) tau\textasciicircum{}-4 at the approach -- the matter-dominated Friedmann value. The radial tidal component likewise -> -4/(9 tau\textasciicircum{}2), likewise M-free. ** SO WHAT DISTINGUISHES THE LOCUS IS THE UNIVERSALITY OF THE APPROACH, NOT THE SIZE OF THE CURVATURE: ** every progenitor spends its last proper interval in the same curvature history, so nothing of the parent survives in the approach. The classification's axis is curvature AT the locus; this is a second axis, and the corpus's scale-freeness is a statement on it. ** DOWNSTREAM: the register's 'well-posed (FINITE-CURVATURE branch point r=0)' is wrong as written -- the curvature diverges and P2's own classification says r=0 is the INFINITE-curvature species at which worldlines TERMINATE. ** So there is no worldline crossing to compute; what continues is the ANALYTIC CONTINUATION. The well-posedness survives on the universal approach and the analyticity rather than on a finiteness that does not hold. ** AND THAT IS WHY THE FIELD HALF HAS SOMETHING TO ACT ON AND THE WORLDLINE HALF DOES NOT: ** the segment has ZERO cosmic duration (Re tautilde does not advance) and FINITE conformal length (\textbar{}Delta eta\textbar{} = 3.32 alpha).). 
\emph{Computes:} K along the bead asserted M-independent; lim tau\textasciicircum{}4 K asserted 64/27; the tidal component asserted M-independent with lim tau\textasciicircum{}2 R = -4/9; \textbar{}Delta eta\textbar{} asserted 3.32
 \emph{Bound:} ** THIS IS A UNIVERSALITY CLAIM ABOUT THE APPROACH, NOT ABOUT THE POINT. ** At fixed r the curvature carries M and diverges; nothing here makes the locus regular or makes worldlines cross it. It adds an axis to the classification and corrects a register description; it overturns no result
\ \emph{Run:} \texttt{python3 receipts/P02\_janzen\_circle/P02\_the\_approach\_is\_mass\_free.py}
\item[\label{rcpt:I8_the_circle_is_a_phase_portrait_and_the_critical_points_are_turning_points}\texttt{I8\_the\_circle\_is\_a\_phase\_portrait\_and\_the\_critical\_points\_are\_turning\_points}]\hfill\textsf{[OK]}\\
\textit{sec:cycloid} \ (I8 --- THE GEOMETRIC CIRCLE IS A PHASE PORTRAIT. Integrable-systems field bake, probe I8. P02 scores x0 on every term in this field's vocabulary and is its founding example. Establishes that r'' = -(r-M) is the harmonic oscillator; that P02's locus (r-M)\textasciicircum{}2 + s\textasciicircum{}2 = M\textasciicircum{}2 IS its phase-space orbit with s = dr/dz; that the conserved energy is M\textasciicircum{}2/2; and that the two critical points are the orbit's two turning points, of which a one-degree-of-freedom periodic orbit has exactly two, interchanged by the time reversal the evenness of r(z) expresses.). 
\emph{Computes:} runs clean; nine verdicts plus an adversarial perturbed curve that fails the orbit identity
 \emph{Bound:} derived here from P02's own eq:cycloid; no parameter pinned
\ \emph{Run:} \texttt{python3 receipts/P02\_janzen\_circle/I8\_the\_circle\_is\_a\_phase\_portrait\_and\_the\_critical\_points\_are\_turning\_points.py}
\item[\label{rcpt:H1_the_FRW_case_is_the_same_cycloid}\texttt{H1\_the\_FRW\_case\_is\_the\_same\_cycloid}]\hfill\textsf{[OK]}\\
\textit{sec:cycloid} \ (the entry-point fronts FRW third is not an analogy: closed FRW dust IS P2s cycloid, exactly, and with the same critical-point character). 
\emph{Computes:} shows the two parametric forms identical under z = pi - eta and A = 2M, both quadratic at the singular end, against P3s own statement of what the argument turns on (11 checks)
 \emph{Bound:} NOT-A-PAPER-CLAIM --- narrows an entry-point site from three cases to two
\ \emph{Run:} \texttt{python3 receipts/L210\_entry\_points/H1\_the\_FRW\_case\_is\_the\_same\_cycloid.py}
\end{description}\endgroup
